\documentclass{exam}

\newif\ifA
\newif\ifkey
\newif\ifprintselected   % required by exam_config's \ansbox

%%%%%%%%%%%%%%%%%%%%%%%%%%%%%%%%%%%%%%%%%%%%%%%%%%%%%%
%%% SET THESE IF VARIABLES
%%%%%%%%%%%%%%%%%%%%%%%%%%%%%%%%%%%%%%%%%%%%%%%%%%%%%%
% Lab name -- labs are named by the lesson they belong to, never by a
% sequential lab number (Neil, 2026-10-12): this is the L17 Lab.
\def\labLesson{17}

% is this the key or not?
\ifdefined\iskey
	\ifnum\iskey=1
		\keytrue
	\else
		\keyfalse
	\fi
\else
	%% Pick one manually
	\keyfalse % if this is NOT the key
%	\keytrue     % if this IS the key
\fi

% set up the header and footer - change for every term
\ifdefined\thisTermIn
	\def\thisterm{\thisTermIn}
\else
	\def\thisterm{Spring 2027}
\fi
%
\def\thisclass{ECE 448}

%%%%%%%%%%%%%%%%%%%%%%%%%%%%%%%%%%%%%%%%%%%%%%%%%%%%%%
%%% END SET THESE IF VARIABLES
%%%%%%%%%%%%%%%%%%%%%%%%%%%%%%%%%%%%%%%%%%%%%%%%%%%%%%
% packages and shortcuts
\input{myPackages_gr}
\usepackage[hidelinks, linkcolor=red]{hyperref}
\input{myShortcuts}
\setlength{\parindent}{0pt}
\setlength{\parskip}{6pt}

%% exam class customization
\input{exam_config}
% ECE 448 document fonts.
%
% Body text is Barlow (SIL OFL; the faces and the license are in
% latex/fonts/Barlow/). Barlow is not in TeX Live, so it is loaded as an
% OpenType face through fontspec -- which means the document must be built
% with lualatex, not pdflatex. build_practice.sh and build_lab.sh both do.
%
% Compiling with pdflatex anyway is not an error: the guard below falls back
% to the LaTeX default so the source still compiles under any engine and can
% be syntax-checked. The resulting PDF just is not the one we ship.
%
% Paths are relative to latex/, which is where both build scripts run
% (each does cd "$(dirname "$0")"). Loading by path rather than by installed
% family name means no system font install and no fontconfig, so a fresh
% clone builds a byte-comparable PDF anywhere.
%
% MATH IS DELIBERATELY LEFT IN COMPUTER MODERN. Barlow has no math companion,
% and CM math under sans body text is already the house pairing: the decks set
% words in Source Sans Pro and let MathJax render symbols in its CM-derived
% face. Moving math to a sans as well would mean unicode-math plus a sans math
% font (Fira Math being the plausible one) and re-checking every equation in
% the course for symbol coverage. That is a separate decision, not a side
% effect of this one.

\usepackage{iftex}
\ifLuaTeX
	\usepackage{fontspec}
	\defaultfontfeatures{Ligatures=TeX}
	\setmainfont{Barlow}[
		Path           = fonts/Barlow/,
		Extension      = .ttf,
		UprightFont    = *-Regular,
		ItalicFont     = *-Italic,
		BoldFont       = *-Bold,
		BoldItalicFont = *-BoldItalic,
	]
	\setsansfont{Barlow}[
		Path           = fonts/Barlow/,
		Extension      = .ttf,
		UprightFont    = *-Regular,
		ItalicFont     = *-Italic,
		BoldFont       = *-SemiBold,
		BoldItalicFont = *-BoldItalic,
	]
	% microtype earns its place here specifically because of Barlow. Its
	% glyphs are narrower than Computer Modern's, so justified prose sets
	% tighter and TeX has to squeeze interword space to fit; two paragraphs
	% in the L02/L03 sets went overfull on the first Barlow build. Protrusion
	% and font expansion absorb that -- both are lualatex features, so this
	% belongs inside the guard with everything else.
	\usepackage{microtype}

	% Mono stays Latin Modern: Barlow ships no monospaced face, and verbatim
	% and code need one that actually is fixed-pitch.
\fi
   % Barlow body text (needs lualatex)
\printselectedfalse      % full worked solutions in the key, not answers-only
\CorrectChoiceEmphasis{\color{red}}
\SolutionEmphasis{\color{red}}
\renewcommand\questionlabel{\textbf{\thequestion.}}
\renewcommand{\thepartno}{\alph{partno}}
\renewcommand{\partlabel}{(\thepartno)}

% Fillable table cell: blank in the student copy, the expected value in red on
% the KEY. Fixed-width centered columns keep the blank cells writable.
\newcolumntype{K}[1]{>{\centering\arraybackslash}p{#1}}
\newcommand{\kv}[1]{\ifkey{\color{red}#1}\fi\rule[-0.10in]{0pt}{0.34in}}

% Answers are not in this public file. Every solution, \kv value, and \ansbox
% value is \keysol{<id>}, defined in the private ECE448-faculty repo at
% lab-keys/src/ECE448_Lab_L17_RFReverseEngineering_solutions.tex; build_lab.sh puts that
% folder on TEXINPUTS (KEYDIR) for the KEY. A KEY build without it stops.
\long\def\defkeysol#1#2{\expandafter\long\expandafter\def\csname keysol@#1\endcsname{#2}}
\def\keysol#1{\ifcsname keysol@#1\endcsname\csname keysol@#1\endcsname\fi}
\ifkey
	\IfFileExists{ECE448_Lab_L17_RFReverseEngineering_solutions.tex}%
		{\input{ECE448_Lab_L17_RFReverseEngineering_solutions}}%
		{\errmessage{KEY build needs ECE448_Lab_L17_RFReverseEngineering_solutions.tex from ECE448-faculty/lab-keys/src (set KEYDIR)}}
\fi

% print the answers if this is the key
\ifkey
	\printanswers % if this IS the key
\else
	\noprintanswers  %if this is NOT the key
\fi

% print the header and footer
% need to print key on top if this is the key
\ifkey
	\firstpageheader{}{\color{red}{ KEY}}{}
	\runningheader{\thisclass\hspace{1pt} -- \thisterm }{L\labLesson\ Lab\hspace{1pt}\color{red}{ KEY}}{\thepage\hspace{1pt} of\hspace{1pt} \numpages}
	\firstpagefooter{}{}{}
	\runningfooter{}{}{}
\else
	\firstpageheader{}{}{}
	\runningheader{\thisclass \hspace{1pt} -- \thisterm }{L\labLesson\ Lab\hspace{1pt}}{\thepage\hspace{1pt} of\hspace{1pt} \numpages}
	\firstpagefooter{}{}{}
	\runningfooter{}{}{}
\fi

\runningheadrule


\begin{document}
\pagestyle{headandfoot}   % exam class

%%%%%%%%%%%%%%%%%%%%%%%%%%%%%%%%%%%%%%%%%%%%%%%%%%%%%%
% title block
%%%%%%%%%%%%%%%%%%%%%%%%%%%%%%%%%%%%%%%%%%%%%%%%%%%%%%

\begin{center}
USAF ACADEMY \\
DEPARTMENT OF ELECTRICAL AND COMPUTER ENGINEERING \\
LESSON \labLesson\ LAB - RF REVERSE ENGINEERING \\
\textbf{Lab Packet} \\
\thisterm \\[4mm]
\end{center}

\vspace{-4pt}
\textbf{Name:} \rule[-2pt]{2.6in}{0.4pt} \hfill
\textbf{Section:} \rule[-2pt]{0.9in}{0.4pt} \hfill
\textbf{Date:} \rule[-2pt]{1.1in}{0.4pt}

\vspace{4pt}

\fullwidth{\textbf{Learning Objective 6.6: } I can perform RF reverse engineering on
an unknown signal: identify its modulation, symbol rate, and payload structure.}

\vspace{2pt}

This packet is the L\labLesson\ Lab, in three parts, and you turn it in once, complete,
before Lesson 18. Do \textbf{Part 1} on your own before the lab period: read the
recording's meta file, check \texttt{L17-chips.grc} on the lesson's recording, and work out the chip settings.
\textbf{Part 2} is what you find, build, and observe, recorded in class: Questions 4 to
7 on your own laptop, Question 8 at a PlutoSDR station with your team, and Question 9
with a real remote of your choice. Do
\textbf{Part 3} on your own after the lab: explain, in words, the evidence for each
answer.

\textbf{Grading.} You submit the packet once, we grade it against Objective 6.6 on
the course's 1 / 0 scale, and you may revise it under course policy. Paste screenshots
where a question asks for one.

\textbf{The device.} The recording \texttt{L17-lab-unknown.sigmf-data}, with its
\texttt{.sigmf-meta} file, is on the L17 Lab page. It is a Simulated remote outlet
switch with six buttons, pressed once each in the order the meta file lists. Nothing
else about it is given: finding it out is the lab.

\textbf{Transmitting.} In class we transmit only into a cable, into a shielded
enclosure, or at very low power in an ISM band, and only when the lesson says to.
This device's band is an FCC Part 15 frequency, not a U.S. ISM band, and a signal the
course makes never goes out on the real device's frequency: your encoder transmits at
915 MHz, in the 902 to 928 MHz ISM band. Questions 1 to 7 run on your laptop and transmit nothing. Question 8 transmits from a
PlutoSDR's TX port, through a cable with an inline 30 dB SMA attenuator, to a second
PlutoSDR's RX port, or, if your instructor says so, over the air at 915 MHz at low
power. Connect the cable and attenuator before you run the flowgraph, and connect an
antenna only when your instructor says so. Never set TX attenuation below 30 dB or change the
slider's range; a Pluto's input is rated to only a few dBm. We have compiled the
PlutoSDR flowgraphs but not run them on a radio, so Question 8 is their first run on
hardware. Question 9 only receives, which is always allowed.

\begin{questions}

\fullwidth{\large\textbf{Part 1: Pre-Lab} \normalsize(individual, before the lab period)}

%%%%%%%%%%%%%%%%%%%%%%%%%%%%%%%%%%%%%%%%%%%%%%%%%%%
\question \textbf{The meta file.} Open \texttt{L17-lab-unknown.sigmf-meta} in a text
editor.

\begin{parts}

	\part Record the recording's datatype, sample rate, and center frequency.

	\begin{center}
	\renewcommand{\arraystretch}{1.25}
	\begin{tabular}{| l | K{1.3in} | K{1.3in} | K{1.5in} |}
	\hline
	& Datatype & Sample rate & Center frequency \\\hline
	\texttt{L17-lab-unknown} & \kv{\keysol{kv1}} & \kv{\keysol{kv2}} & \kv{\keysol{kv3}} \\\hline
	\end{tabular}
	\end{center}

	\part Give the values you set for \texttt{rec}, \texttt{samp\_rate}, and \texttt{fc}
	in \texttt{L17-sigmf-view.grc}, and the Output Type of its SigMF Source.
		\begin{solution}[0.9in]\keysol{sol1}\end{solution}

	\part Paste a screenshot of the viewer running on this recording.
		\begin{solution}[1.6in]\keysol{sol2}\end{solution}

\end{parts}

%%%%%%%%%%%%%%%%%%%%%%%%%%%%%%%%%%%%%%%%%%%%%%%%%%%
\newpage
\question \textbf{Reading chips in GNU Radio.} Open the lesson's recording,
\texttt{L17-fan-remote.sigmf-data}, in \texttt{L17-chips.grc}, with \texttt{samp\_rate}
200000, \texttt{fc} 350e6, and the carrier offset at 40000 Hz.

\begin{parts}

	\begin{minipage}[t]{0.58\linewidth}
		\part Measure the shortest pulse on the envelope, and set the chip length to it.
		Record the chip length in samples and in microseconds, and the first line of chips
		the Chip Reader prints, with its pause.
		\begin{solution}[0.9in]\keysol{sol3}\end{solution}
	\end{minipage}\hfill%
	\begin{minipage}[t]{0.37\linewidth}
		\raggedleft \vspace{12pt}
		\ansbox[1.4in]{\keysol{ans1}}
	\end{minipage}

	\part Paste a screenshot of \texttt{L17-chips.grc} running on the recording.
		\begin{solution}[1.6in]\keysol{sol4}\end{solution}

\end{parts}

%%%%%%%%%%%%%%%%%%%%%%%%%%%%%%%%%%%%%%%%%%%%%%%%%%%
\question \textbf{Samples per chip.} The Chip Reader's chip length and your encoder's
\texttt{chip\_len} are both a chip's length in samples. Fill in the table for each chip
length at this recording's sample rate, from Question 1(a), and at the PlutoSDR
flowgraph's 1 MS/s.

	\begin{center}
	\renewcommand{\arraystretch}{1.25}
	\begin{tabular}{| l | K{0.8in} | K{0.8in} | K{0.8in} | K{0.8in} |}
	\hline
	Chip length & 240 $\mu$s & 300 $\mu$s & 600 $\mu$s & 1 ms \\\hline
	At the recording's rate & \kv{\keysol{kv4}} & \kv{\keysol{kv5}} & \kv{\keysol{kv6}} & \kv{\keysol{kv7}} \\\hline
	At 1 MS/s & \kv{\keysol{kv8}} & \kv{\keysol{kv9}} & \kv{\keysol{kv10}} & \kv{\keysol{kv11}} \\\hline
	\end{tabular}
	\end{center}

		\begin{solution}[0.7in]\keysol{sol5}\end{solution}

%%%%%%%%%%%%%%%%%%%%%%%%%%%%%%%%%%%%%%%%%%%%%%%%%%%
\newpage
\fullwidth{\large\textbf{Part 2: Measurements} \normalsize(in class)}

\question \textbf{Step 1: the carrier.} Follow the lab page's steps for
\texttt{L17-sigmf-view.grc}.

\begin{parts}

	\part List every signal on the waterfall, its frequency, and whether it bursts or
	stays on. Say which one is the device, and why.
		\begin{solution}[1.1in]\keysol{sol6}\end{solution}

	\begin{minipage}[t]{0.58\linewidth}
		\part Record the carrier offset: the center of the bursts on the Frequency Sink or
		waterfall, checked as the midpoint of the half-height offsets on the envelope.
		Paste a screenshot of the window at that offset.
		\begin{solution}[1.6in]\keysol{sol7}\end{solution}
	\end{minipage}\hfill%
	\begin{minipage}[t]{0.37\linewidth}
		\raggedleft \vspace{12pt}
		\ansbox[1.3in]{\keysol{ans2}}
	\end{minipage}

	\begin{minipage}[t]{0.58\linewidth}
		\part Record the device's carrier frequency.
		\begin{solution}[0.6in]\keysol{sol8}\end{solution}
	\end{minipage}\hfill%
	\begin{minipage}[t]{0.37\linewidth}
		\raggedleft \vspace{12pt}
		\ansbox[1.3in]{\keysol{ans3}}
	\end{minipage}

\end{parts}

%%%%%%%%%%%%%%%%%%%%%%%%%%%%%%%%%%%%%%%%%%%%%%%%%%%
\newpage
\question \textbf{Steps 2 and 3: modulation and encoding.} Follow the lab page's steps in
\texttt{L17-chips.grc}.

\begin{parts}

	\begin{minipage}[t]{0.58\linewidth}
		\part Record the shortest pulse you measured on the envelope, in samples and in
		microseconds, and the chip length you set.
		\begin{solution}[0.7in]\keysol{sol9}\end{solution}
	\end{minipage}\hfill%
	\begin{minipage}[t]{0.37\linewidth}
		\raggedleft \vspace{12pt}
		\ansbox[1.3in]{\keysol{ans4}}
	\end{minipage}

	\part Record the shortest gap you measured, in samples and in microseconds. Say
	whether it fits your chip length, and how you checked that the Chip Reader's lines
	repeat from burst to burst.
		\begin{solution}[1.2in]\keysol{sol10}\end{solution}

	\part Name the modulation and the encoding. Give the evidence for each from the
	envelope and the lines of chips, and describe any structure in each burst besides
	the data bits, if there is any.
		\begin{solution}[1.4in]\keysol{sol11}\end{solution}

	\part Find the baud rate, in chips per second, and the bit rate, in data bits per
	second.

	\begin{center}
	\renewcommand{\arraystretch}{1.25}
	\begin{tabular}{| l | K{1.2in} | K{1.2in} |}
	\hline
	& Baud rate & Bit rate \\\hline
	The device & \kv{\keysol{kv12}} & \kv{\keysol{kv13}} \\\hline
	\end{tabular}
	\end{center}

		\begin{solution}[0.7in]\keysol{sol12}\end{solution}

\end{parts}

%%%%%%%%%%%%%%%%%%%%%%%%%%%%%%%%%%%%%%%%%%%%%%%%%%%
\newpage
\question \textbf{Step 3: the fields.}

\begin{parts}

	\part Decode one burst of each button and write its data bits in hex.

	\begin{center}
	\renewcommand{\arraystretch}{1.25}
	\begin{tabular}{| l | K{1.4in} || l | K{1.4in} |}
	\hline
	1 On & \kv{\keysol{kv14}} & 2 Off & \kv{\keysol{kv17}} \\\hline
	1 Off & \kv{\keysol{kv15}} & 3 On & \kv{\keysol{kv18}} \\\hline
	2 On & \kv{\keysol{kv16}} & 3 Off & \kv{\keysol{kv19}} \\\hline
	\end{tabular}
	\end{center}

	\part Mark the fields: which bits never change, which change with the button, and
	any field that follows neither. If you find such a field, give the rule it follows
	and check it on two buttons.
		\begin{solution}[1.6in]\keysol{sol13}\end{solution}

	\begin{minipage}[t]{0.58\linewidth}
		\part How many times does one press send its frame, and how long is the pause
		between frames?
		\begin{solution}[0.8in]\keysol{sol14}\end{solution}
	\end{minipage}\hfill%
	\begin{minipage}[t]{0.37\linewidth}
		\raggedleft \vspace{12pt}
		\ansbox[1.3in]{\keysol{ans5}}
	\end{minipage}

\end{parts}

%%%%%%%%%%%%%%%%%%%%%%%%%%%%%%%%%%%%%%%%%%%%%%%%%%%
\newpage
\question \textbf{Step 4: build the encoder.} Start GRC with an empty canvas. Do not
open \texttt{L17-fan-tx-sim.grc}: build the encoder yourself as the lab page lists,
with your settings from Questions 4 to 6.

\begin{parts}

	\part Paste a screenshot of your flowgraph's canvas, and the lines of chips
	\texttt{L17-chips.grc} prints for \texttt{my-unknown.sigmf-data} and for the
	original recording, for the same button.
		\begin{solution}[1.8in]\keysol{sol15}\end{solution}

	\part Do the same for a second button, and say whether every chip matches.
		\begin{solution}[1.2in]\keysol{sol16}\end{solution}

	\begin{minipage}[t]{0.58\linewidth}
		\part Record the noise amplitude at which the Chip Reader's lines of chips from your
		recording first stop matching the original, and describe what goes wrong first.
		\begin{solution}[1.0in]\keysol{sol17}\end{solution}
	\end{minipage}\hfill%
	\begin{minipage}[t]{0.37\linewidth}
		\raggedleft \vspace{12pt}
		\ansbox[1.3in]{\keysol{ans6}}
	\end{minipage}

\end{parts}

%%%%%%%%%%%%%%%%%%%%%%%%%%%%%%%%%%%%%%%%%%%%%%%%%%%
\newpage
\question \textbf{Step 5: transmit into a cable.} At a PlutoSDR station, with your team,
follow the lab page's steps: connect the first Pluto's TX port to the second Pluto's RX port
through the 30 dB attenuator, set \texttt{fc} to 915e6 and run \texttt{L17-capture.grc}
on the second Pluto, and run your Pluto encoder at 915 MHz. When you finish, stop both flowgraphs and unplug the radios before
you disconnect the cable.

\begin{parts}

	\part Paste a screenshot of both windows with your signal on the capture
	flowgraph's waterfall.
		\begin{solution}[1.5in]\keysol{sol18}\end{solution}

	\part Open \texttt{capture.sigmf-data} in \texttt{L17-chips.grc} at the sample rate
	and center frequency in its meta file, 1 MS/s and 915 MHz, and compare its lines of
	chips with the original recording's for the same button.
		\begin{solution}[1.0in]\keysol{sol19}\end{solution}

	\part Record whether the chips still match at each TX attenuation, and explain in words,
	from what you saw on the waterfall and in \texttt{L17-chips.grc}, what limits the link as the attenuation
	rises.

	\begin{center}
	\renewcommand{\arraystretch}{1.25}
	\begin{tabular}{| l | K{0.9in} | K{0.9in} | K{0.9in} | K{0.9in} |}
	\hline
	TX attenuation (dB) & 60 & 70 & 80 & 89 \\\hline
	Chips match? & & & & \\\hline
	\end{tabular}
	\end{center}

		\begin{solution}[1.0in]\keysol{sol20}\end{solution}

\end{parts}

%%%%%%%%%%%%%%%%%%%%%%%%%%%%%%%%%%%%%%%%%%%%%%%%%%%
\newpage
\question \textbf{A real remote (receive only).} Follow the lab page's steps with a real
key fob, doorbell, or sensor of your choice. Answers depend on the device.

\begin{parts}

	\part Name the device, and record its carrier frequency and how you found it:
	the spectrum of \texttt{L17-explore-pluto.grc}. Paste a screenshot.
		\begin{solution}[1.4in]\keysol{sol24}\end{solution}

	\part Record the modulation, the chip length, and the encoding, with the evidence
	from \texttt{L17-chips.grc} for each.
		\begin{solution}[1.2in]\keysol{sol25}\end{solution}

	\part Compare the recordings of two buttons, and of two presses of the same button.
	Say which bits change, and whether a press repeats the same frame each time.
		\begin{solution}[1.2in]\keysol{sol26}\end{solution}

\end{parts}

%%%%%%%%%%%%%%%%%%%%%%%%%%%%%%%%%%%%%%%%%%%%%%%%%%%
\newpage
\fullwidth{\large\textbf{Part 3: Writeup} \normalsize(individual, after the lab)}

\question \textbf{Analysis.}

\begin{parts}

	\part For each of the carrier, the modulation, the encoding, and each field, give
	the observation that settled your answer: what you saw on the waterfall, the
	envelope, or the lines of chips, and which comparison between buttons identified
	each field.
		\begin{solution}[1.8in]\keysol{sol21}\end{solution}

	\part From your Question 6 table and your noise run in Question 7(c), explain how a
	receiver of this device could tell a frame damaged by noise from a good one, and
	what kinds of damage it could miss.
		\begin{solution}[1.5in]\keysol{sol22}\end{solution}

	\part Explain why your encoder from Question 7 would operate a device that matched
	this recording, and
	describe in words how a rolling code, as in a car key fob, would stop a recording of
	one press from being reused.
		\begin{solution}[1.5in]\keysol{sol23}\end{solution}

\end{parts}

\end{questions}

\vspace{1.0em}
\noindent\textbf{Documentation:} \rule[-2pt]{0.6\linewidth}{0.4pt}

\end{document}
